\chapter{Приложение А. Дополнительные синтетические эксперименты и проверки устойчивости}
\label{app:prilozhenie-a-dopolnitelnye-synthetic-eksperimenty-i-proverki-ustoichivosti}

В основное тело диплома включены только те результаты синтетического эксперимента, которые непосредственно поддерживают центральный тезис работы. Остальные материалы вынесены в приложение.

В приложение А входят:

\begin{enumerate}
\def\labelenumi{\arabic{enumi}.}
\tightlist
\item
  Полные серии однотипных графиков по базовой, запаздывающей и смешанной моделям, включая дополнительные режимы, фазовые портреты и прогнозные продолжения.
\item
  Дополнительные переборы параметров и сценарии по \(q\) и \(\gamma\) смешанной модели.
\item
  Развёрнутые robustness checks по длине окна, сравнению \(B\) и \(\beta\), критерию интерпретируемости, регуляризации, PCA и настройке \(p\)-порога.
\item
  Полные таблицы метрик по методам, окнам и режимам, которые были использованы при сжатии раздела 2.4 до основной бакалаврской версии.
\end{enumerate}

Электронные материалы приложения А соответствуют следующим наборам файлов:

\begin{itemize}
\tightlist
\item
  \texttt{outputs/planA} --- direct vs regression, шумовой sweep и диагностическая полезность;
\item
  \texttt{outputs/planB} --- влияние длины окна и отношение \(W/T\);
\item
  \texttt{outputs/planC} --- сравнение коэффициентов \(B\) и \(\beta\);
\item
  \texttt{outputs/planD} --- ENTER, STEPWISE и регуляризованные методы;
\item
  \texttt{outputs/planE} --- чувствительность к \(p\)-порогу;
\item
  \texttt{outputs/planF} --- PCA и ортогонализация лагового пространства;
\item
  \texttt{outputs/planG} --- режимная предклассификация и выбор контура;
\item
  \texttt{outputs/planH} --- дисперсия, кондиционность и риск вырождения окна.
\end{itemize}
